\section{Reproducibility specification}
\label{app:repro}

This appendix specifies each family's selected vision layer, the consumed-token mask, the projection and probe, the control constructions, the cohort roles, and the protocol record at the registered protocol's level of detail. The released code implements each component.

\subsection{Consumed sites per family}
\label{app:repro-sites}

The selected vision layer is the last visual block that feeds the consumer. The connector site is the consumer's output that enters the language model. We identify the block from each checkpoint's own configuration when loading the checkpoint. Each block's preflight record stores the identified block, its hidden dimension, its consumer, and any branch that bypasses it.

\begin{itemize}[leftmargin=*,itemsep=1pt,topsep=2pt]
\item Qwen2.5-VL, Lingshu, and Qwen3-VL: the merger reads the last block of the vision transformer, which has 32 blocks in Qwen2.5-VL and Lingshu, 24 in Qwen3-VL 4B, and 27 in Qwen3-VL 8B and 32B. A Qwen2.5-VL or Lingshu image contributes 576 patch tokens, which the merger joins in $2\times2$ groups into 144; a Qwen3-VL image contributes at most 1{,}024. The Qwen3-VL DeepStack mergers, which read earlier blocks, are not written.
\item Gemma~3 and MedGemma: the projector reads the last of the 27 SigLIP encoder layers through the tower's final layer normalisation, and consumes all 4{,}096 patch tokens of the $896\times896$ image.
\item InternVL3.5: the projector reads the last InternViT encoder layer, of 24 at 8B and 14B and of 45 at 38B; the tower's final normalisation is an identity in the released configuration. An image contributes 1{,}024 patch tokens from a single $448\times448$ tile.
\item LLaVA-1.5 and LLaVA-Med: the projector reads the penultimate of the 24 CLIP encoder layers, and the final layer is computed and discarded. An image contributes 576 patch tokens.
\item Llama~3.2 Vision: the projector input concatenates the last of the eight global-encoder layers with five intermediate layers of the local encoder, which bypass the selected vision layer. An image contributes 1{,}601 tokens, the CLS token and 1{,}600 patches, from its one real $560\times560$ tile.
\end{itemize}
Table~\ref{tab:cf-models} lists the consumed-token count of every checkpoint.

\subsection{The consumed-token mask}
\label{app:repro-mask}

The write in Equation~\ref{eq:intervention} changes only the consumed tokens at the positions the consumer reads at the site. We derive these positions for each batch from that batch's processor outputs:

\begin{itemize}[leftmargin=*,itemsep=1pt,topsep=2pt]
\item Qwen2.5-VL, Lingshu, and Qwen3-VL: the tokens of all images in a batch form one sequence in which every image occupies one contiguous span, whose bounds follow from the image's patch grid. At the connector site the span is a quarter as long, one merged token per $2\times2$ patch group.
\item Gemma~3 and MedGemma: all 4{,}096 patch tokens; the tower produces no CLS token. The number of written tokens must equal the number of image placeholder tokens in the language-model input.
\item InternVL3.5: the patch tokens of the single $448\times448$ tile. The CLS token, which the consumer drops, is excluded. The number of written tokens must equal the number of image placeholder tokens.
\item LLaVA-1.5 and LLaVA-Med: the 576 patch tokens. The CLS token, which the checkpoint's feature selection drops, is excluded.
\item Llama~3.2 Vision: the 1{,}601 tokens, CLS and 1{,}600 patches, of the one real $560\times560$ tile. The seven padding tokens of each tile and the three padding tiles are excluded, in agreement with the model's own aspect-ratio and cross-attention masks.
\end{itemize}

Scoring stops with an error and no fallback when an image span does not exactly cover the site, a mask length differs from the token count at the site, an image has no consumed token, or one forward pass reaches the site more than once.

Preflight check (C) verifies the mask. It captures the site tensor on a clean pass and a pass steered at $\alpha=+0.25$ along the first concept direction. It reports the maximum absolute difference outside the mask and the maximum absolute changes in the consumer's output and answer logits. Across all \cfBlocks{} blocks, the change outside the consumed tokens is $0$ at both sites. The consumer's output and answer logits both change. The same records show a bitwise-identical repeated clean pass (check A) and a bitwise-identical $\alpha=0$ pass (check B) for every block.

\subsection{Projection, lift, scaler, and probe}
\label{app:repro-probe}

We construct the projection $R\in\R^{D\times512}$ by drawing standard normal variates once from a PCG64 generator seeded with $0$, dividing by $\sqrt{512}$, and storing the result in float32. Its entries are independent $\mathcal{N}(0,1/512)$ variates. The seed is $0$ for every checkpoint, dataset, and site. $D$ is the family's hidden dimension. We neither normalise nor orthogonalise the columns of $R$. The Gram matrix $R^{\top}R/D$ has full rank 512 in all \cfAltdirdBlocks{} blocks that carry the displacement module, with eigenvalues $\cfAltdirdEigMin$--$\cfAltdirdEigMax$ and condition number \cfAltdirdCondMin{}--\cfAltdirdCondMax{}.

Projecting the token-mean read vector $\hvec_i$ gives features $R^{\top}\hvec_i$. A featurewise standardiser fitted only on training rows supplies the mean $\boldsymbol\mu$ and scale $\boldsymbol s$. Every row is standardised as $(R^{\top}\hvec-\boldsymbol\mu)\oslash\max(\boldsymbol s,10^{-8})$, where $\oslash$ and $\max$ act elementwise. Each concept probe is an $\ell_2$-regularised logistic regression with inverse regularisation strength $C=1$. We fit each probe by L-BFGS for at most 2{,}000 iterations with unweighted classes and seed $0$, using training rows with a known label for that concept. For each of the two refit seeds, we resample training units with replacement, retaining all rows of each sampled unit at its multiplicity. We then refit the scaler and six probes using the same projection. The random family, shams, and control probes remain at seed 0.

\subsection{Control constructions}
\label{app:repro-controls}

\paragraph{Random directions.} We draw one $119\times D$ standard-normal matrix using a PCG64 generator seeded with $0$ and normalise each row to unit $\ell_2$ norm. Every question in the block uses the same 119 directions.

\paragraph{Coordinate-permutation sham.} Immediately after the random draw, the same generator produces one permutation of $\{0,\dots,D{-}1\}$ per concept, in protocol concept order. We form concept $c$'s sham by rearranging the coordinates of $\hat\wvec_c$ with concept $c$'s own permutation. This preserves the direction's coordinate multiset and unit norm. We write each question with its own sham.

\paragraph{Matched random-label probes.} Each row has a type: view $\times$ sex $\times$ age decade on the chest sets, and aspect-ratio $\times$ area buckets on COCO. For control seed $k\in\{0,\dots,19\}$, we traverse the types in sorted order and use a PCG64 generator seeded with $k$ to draw one label per type uniformly from $\{0,1\}$. If all types receive the same label, we flip the first. Every row of a type receives that type's label. We give each control probe the concept probe's capacity and nuisance structure by fitting it on the same training rows and known-label subset with the same hyperparameters. A block fits control probes when its training rows span at least eight types. Selectivity is Equation~\ref{eq:selectivity}. A cell is graded for readability when the real AUROC is finite and at least ten of the twenty control AUROCs are finite.

\subsection{Cohort roles}
\label{app:repro-cohorts}

Every chest image is a frontal study. One role holds more than one row per unit: the NIH training role, which takes every study of each training patient, \cfCohortNihTrainRows{} studies from \cfCohortNihTrainUnits{} patients, \cfCohortNihTrainMulti{} of whom contribute more than one and one of whom contributes \cfCohortNihTrainMax{}. Every other role of every dataset holds one row per unit, and no unit appears in two roles of a dataset. We use the calibration role only to determine readability and answerability. We neither fit probes nor score writes on it. Every write uses the test role.

\FloatBarrier

\subsection{Protocol amendments}
\label{app:repro-protocol}

Every analysis the paper reports beyond the \cfCoveragePlannedWord{} planned modules was added to the protocol file as a dated amendment after the first wave, and each amendment precedes the first result of the module it registers. Each amendment entry names its own rows, references, and decision rule before it names a result.


\FloatBarrier

\paragraph{The fine-grained category selection rule.} The rule that picks the six fine-grained COCO categories is stated in the protocol file in full, and it reads on the labels alone. A category is eligible when its training positives are at least as many as the weakest of the six existing concepts, when it has at least ten positives among the 400 calibration rows and at least ten among the 600 test rows, and when its mean relative instance area is below the median of the six existing concepts. Eligible categories are ranked by calibration positives in descending order, ties by mean relative instance area in ascending order, and remaining ties by category identifier; the first six are taken. The rule uses the label counts and the instance areas of the fixed cohorts and nothing else: no probe, no answer, and no write enters it. The protocol file carries the rule under the date of its amendment, and the module's first result is later than that date, so the six categories and the two matching windows were fixed before any of them was scored.
